数学公式在实数上成立，程序却只在有限集合中运算。数值分析的第一课是学会拆分三种来源：输入或问题本身是否敏感，离散近似省略了什么，算法怎样传播每一步舍入。三件事常以同一症状出现——计算结果跑偏——但纠正手段完全不同。一个小例子：解

\[
\begin{cases} x + y = 2 \\ 1.0001x + y = 2.0001 \end{cases}
\]

用 Gauss 消去，若以 4 位十进制精度运算，算出的 \(y\) 大约是 \(1.0\)，\(x\) 大约是 \(1.0\)——真解是 \(x=1,\;y=1\)，但换个接近奇异的问题，4 位精度就会崩掉。问题出在条件数还是算法？条件数看问题，稳定性看算法，精度看两者乘积。只有分开三件事，才知道该换算法、换精度、换模型，还是承认数据已经不足以决定答案。

建议依次阅读：

\begin{enumerate}
\tightlist
\item
  从实数到浮点数
\item
  误差如何产生与放大
\item
  条件数衡量问题本身的敏感性
\item
  前向误差、后向误差与稳定算法
\end{enumerate}

第一次阅读不必记得 binary64 的每个字段。让每一次出错都能回答三问：输入稍微变一点，正确答案变多少；所用公式省略了哪些项；浮点运算额外丢了多少位。三问分别指向条件性、截断和舍入，分开之后，提高精度这件事的收益才看得清。

前一单元：\hyperref[ch:059]{``信息与学习：哪些数据值得看，哪些表示值得保留''}。

\subsection{本章篇目}

以下篇目按概念依赖排列；若从具体问题进入，也可以先打开最接近当前困惑的一篇，再沿篇内链接回补。

\begin{itemize}
\tightlist
\item \hyperref[sec:08-Numerical-Analysis/01-Floating-Point-and-Error/01]{``从实数到浮点数''}；
\item \hyperref[sec:08-Numerical-Analysis/01-Floating-Point-and-Error/02]{``误差如何产生与放大''}；
\item \hyperref[sec:08-Numerical-Analysis/01-Floating-Point-and-Error/03]{``条件数衡量问题本身的敏感性''}；
\item \hyperref[sec:08-Numerical-Analysis/01-Floating-Point-and-Error/04]{``前向误差、后向误差与稳定算法''}。
\end{itemize}
